% =========================================================================
\section{Homogénea con coeficientes constantes}\label{sec:constantes}
% =========================================================================

Supongamos que $a_0,\dots,a_n\in\R$ son constantes, con $a_n\neq0$. Entonces
$\Lop$ es un \emph{polinomio en el operador} $D$:
\begin{equation}\label{eq:pD}
  \Lop=p(D),\qquad p(\lambda)\defeq a_n\lambda^{n}+a_{n-1}\lambda^{n-1}+\dots+a_1\lambda+a_0 ,
\end{equation}
donde $p$ es el \emph{polinomio característico}. Trabajamos con funciones
de valores complejos y extraemos al final las soluciones reales
\cite[cap.~IV]{TEMARIO}, \cite{ELSGOLTZ}, \cite[cap.~4]{BOYCE}.

% -------------------------------------------------------------------------
\subsection{El ansatz exponencial}\label{sec:ansatz}
% -------------------------------------------------------------------------

\begin{proposicion}[La exponencial es autofunción de $D$]\label{prop:autofuncion}
Para todo $\lambda\in\C$,
\begin{equation}\label{eq:autofuncion}
  \Lop\bigl[\eul^{\lambda x}\bigr]=p(\lambda)\,\eul^{\lambda x}.
\end{equation}
En consecuencia, $\eul^{\lambda x}\in\Ker\Lop$ si y sólo si $\lambda$ es raíz de
la \emph{ecuación característica} $p(\lambda)=0$.
\end{proposicion}

\begin{proof}
$D^{j}\eul^{\lambda x}=\lambda^{j}\eul^{\lambda x}$, luego
$\Lop[\eul^{\lambda x}]=\sum_j a_j\lambda^{j}\eul^{\lambda x}=p(\lambda)\eul^{\lambda x}$,
y $\eul^{\lambda x}\neq0$.
\end{proof}

% -------------------------------------------------------------------------
\subsection{Factorización del operador}\label{sec:factorizacion}
% -------------------------------------------------------------------------
Por el teorema fundamental del álgebra,
\begin{equation}\label{eq:factorizacion-p}
  p(\lambda)=a_n\prod_{k=1}^{r}(\lambda-\lambda_k)^{m_k},
  \qquad \lambda_k\in\C\ \text{distintas},\qquad \sum_{k=1}^{r}m_k=n .
\end{equation}

\begin{proposicion}[Factorización de $\Lop$]\label{prop:factorizacion}
Los operadores $D-\lambda_k$ conmutan entre sí y
\begin{equation}\label{eq:factorizacion-L}
  \Lop=a_n\,(D-\lambda_1)^{m_1}(D-\lambda_2)^{m_2}\cdots(D-\lambda_r)^{m_r},
\end{equation}
con los factores en cualquier orden.
\end{proposicion}

\begin{proof}
Como $D^{i}D^{j}=D^{i+j}$ y las constantes conmutan con $D$, la aplicación
$q(\lambda)\mapsto q(D)$ respeta sumas y productos de polinomios con
coeficientes constantes. Por tanto el producto \eqref{eq:factorizacion-p} se
traslada a \eqref{eq:factorizacion-L}, y $q_1(D)\,q_2(D)=(q_1q_2)(D)=q_2(D)\,q_1(D)$.
\end{proof}

\begin{corolario}[Raíces simples]\label{cor:simples}
Si todas las raíces son simples ($r=n$, $m_k=1$), las exponenciales
$\eul^{\lambda_1x},\dots,\eul^{\lambda_nx}$ forman una base de $\Ker\Lop$, y
\begin{equation}\label{eq:simples}
  y_h(x)=c_1\eul^{\lambda_1x}+c_2\eul^{\lambda_2x}+\dots+c_n\eul^{\lambda_nx}.
\end{equation}
\end{corolario}

\begin{proof}
Son soluciones por la \cref{prop:autofuncion}. En $x_0=0$, la matriz
wronskiana es la de Vandermonde $\bigl(\lambda_k^{\,j}\bigr)$, de determinante
$\prod_{i<k}(\lambda_k-\lambda_i)\neq0$; se concluye por la
\cref{prop:wronskiano}.
\end{proof}

% -------------------------------------------------------------------------
\subsection{Soluciones reales}\label{sec:reales}
% -------------------------------------------------------------------------

\begin{proposicion}[Raíces complejas conjugadas]\label{prop:reales}
Si los $a_j$ son reales y $\lambda=\alpha+\iu\beta$, $\beta\neq0$, es raíz de
$p$, también lo es $\bar\lambda$, con la misma multiplicidad, y el par
$\{\eul^{\lambda x},\eul^{\bar\lambda x}\}$ puede sustituirse por el par
real
\begin{equation}\label{eq:reales}
  \Real\eul^{\lambda x}=\eul^{\alpha x}\cos\beta x,\qquad
  \Imag\eul^{\lambda x}=\eul^{\alpha x}\sin\beta x ,
\end{equation}
que genera el mismo subespacio.
\end{proposicion}

\begin{proof}
Como $p$ tiene coeficientes reales, $p(\bar\lambda)=\overline{p(\lambda)}=0$.
Al ser $\Lop$ real, $\Lop[\Real y]=\Real\Lop[y]$ e igual con la parte
imaginaria, luego \eqref{eq:reales} son soluciones. Por último,
$\eul^{\lambda x}=\eul^{\alpha x}(\cos\beta x+\iu\sin\beta x)$ y
$\eul^{\bar\lambda x}=\eul^{\alpha x}(\cos\beta x-\iu\sin\beta x)$, de
modo que ambos pares generan el mismo plano.
\end{proof}

\begin{ejemplo}[$2y''-y'-y=0$]\label{ej:simple}
$p(\lambda)=2\lambda^{2}-\lambda-1=(2\lambda+1)(\lambda-1)$, de raíces
$\lambda_1=1$ y $\lambda_2=-\tfrac12$. Por el \cref{cor:simples},
\begin{equation}\label{eq:ej-simple}
  y(x)=c_1\,\eul^{x}+c_2\,\eul^{-x/2}.
\end{equation}
\end{ejemplo}

\begin{ejemplo}[$y''+2y'+2y=0$ \cite{TEMARIO}]\label{ej:complejas}
$p(\lambda)=\lambda^{2}+2\lambda+2$ tiene raíces $\lambda=-1\pm\iu$, y por la
\cref{prop:reales}
$y(x)=\eul^{-x}\bigl(k_1\cos x+k_2\sin x\bigr)$.
\end{ejemplo}

% -------------------------------------------------------------------------
\subsection{El oscilador amortiguado}\label{sec:amortiguado}
% -------------------------------------------------------------------------

\begin{ejemplo}[$y''+\nu\omega\,y'+\omega^{2}y=0$, con $\nu\ge0$, $\omega>0$]\label{ej:amortiguado}
La ecuación característica $\lambda^{2}+\nu\omega\lambda+\omega^{2}=0$ da
\begin{equation}\label{eq:raices-oscilador}
  \lambda_\pm=\frac{-\nu\omega\pm\sqrt{\nu^{2}\omega^{2}-4\omega^{2}}}{2}
             =\frac{\omega}{2}\Bigl(-\nu\pm\sqrt{\nu^{2}-4}\Bigr).
\end{equation}
El signo del discriminante $\nu^{2}-4$ decide el comportamiento
(\cref{fig:amortiguado}):
\begin{itemize}
  \item $\nu>2$ (\emph{sobreamortiguado}): dos raíces reales negativas y
        $y=c_1\eul^{\lambda_+x}+c_2\eul^{\lambda_-x}$; decae exponencialmente sin
        oscilar.
  \item $\nu<2$ (\emph{subamortiguado}): raíces complejas conjugadas
        $\lambda_\pm=-\gamma\pm\iu\Omega$, con $\gamma=\nu\omega/2$ y
        $\Omega=\omega\sqrt{1-\nu^{2}/4}$. Por la \cref{prop:reales},
        \begin{equation}\label{eq:subamortiguado}
          y=\eul^{-\gamma x}\bigl(A\cos\Omega x+B\sin\Omega x\bigr)
          \qquad\text{(oscilaciones amortiguadas)}.
        \end{equation}
  \item $\nu=2$ (\emph{crítico}): raíz doble $\lambda=-\omega$. Sólo se
        obtiene una exponencial, $\eul^{-\omega x}$, y falta una segunda
        solución. Este caso frontera motiva la \cref{sec:multiples}.
\end{itemize}
Como $\lambda_+\lambda_-=\omega^{2}$, para $\nu<2$ las raíces recorren la
circunferencia $\abs{\lambda}=\omega$; en $\nu=2$ chocan en $-\omega$ y para
$\nu>2$ se separan sobre el eje real.
\end{ejemplo}

% Figura: lugar de las raíces y respuestas del oscilador amortiguado (omega = 1)
\begin{figure}[ht]
\centering
\begin{tikzpicture}
  % ---------------- (a) Raíces en el plano complejo ----------------
  \begin{axis}[name=raices, oscuro, grid=none, axis equal image,
      width=6.4cm, height=6.4cm,
      xmin=-3.0, xmax=0.7, ymin=-1.55, ymax=1.55,
      xtick={-3,-2,-1,0}, ytick={-1,0,1},
      xlabel={$\Real\lambda$}, ylabel={$\Imag\lambda$},
      title={(a) Raíces $\lambda_\pm(\nu)$, $\omega=1$}]
    \addplot[guia, domain=0:360, samples=100] ({cos(x)},{sin(x)});
    \draw[gris, line width=0.4pt] (axis cs:-3,0) -- (axis cs:0.7,0);
    \draw[gris, line width=0.4pt] (axis cs:0,-1.55) -- (axis cs:0,1.55);
    % 0 <= nu <= 2: arco de la circunferencia unidad (semiplano izquierdo)
    \addplot[violeta, line width=1.3pt, domain=90:180, samples=60] ({cos(x)},{sin(x)});
    \addplot[violeta, line width=1.3pt, domain=180:270, samples=60] ({cos(x)},{sin(x)});
    % nu >= 2: eje real
    \addplot[naranja, line width=1.3pt] coordinates {(-1,0) (-0.38,0)};
    \addplot[naranja, line width=1.3pt] coordinates {(-1,0) (-2.62,0)};
    % ejemplos: nu = 0.5 (sub), 2 (crit), 3 (sobre)
    \addplot[only marks, mark=*, mark size=2.6pt, cian] coordinates {(-0.25,0.968) (-0.25,-0.968)};
    \addplot[only marks, mark=*, mark size=2.8pt, verde] coordinates {(-1,0)};
    \addplot[only marks, mark=*, mark size=2.6pt, magenta] coordinates {(-0.382,0) (-2.618,0)};
    \node[font=\scriptsize, text=gris, anchor=south] at (axis cs:-1.9,1.05) {$\abs{\lambda}=\omega$};
    \node[font=\scriptsize, text=verde, anchor=north east] at (axis cs:-1.15,-0.15) {$\nu=2$};
  \end{axis}
  % ---------------- (b) Respuestas con y(0)=1, y'(0)=0 ----------------
  \begin{axis}[at={($(raices.east)+(1.35cm,0)$)}, anchor=west, oscuro,
      width=8.0cm, height=5.6cm, xmin=0, xmax=14, ymin=-0.8, ymax=1.1,
      xtick={0,2,...,14}, ytick={-0.5,0,0.5,1}, xlabel={$x$},
      title={(b) Respuestas con $y(0)=1$, $y'(0)=0$},
      samples=300, domain=0:14, legend pos=north east]
    \addplot[c1] {exp(-0.25*x)*(cos(deg(0.968246*x))+0.258199*sin(deg(0.968246*x)))};
    \addplot[c4] {(1+x)*exp(-x)};
    \addplot[c2] {(-2.618034*exp(-0.381966*x)+0.381966*exp(-2.618034*x))/(-2.236068)};
    \legend{$\nu=0{,}5$ (sub), $\nu=2$ (crítico), $\nu=3$ (sobre)}
  \end{axis}
\end{tikzpicture}
\caption{Oscilador amortiguado con $\omega=1$. (a) Al crecer $\nu$ desde
$0$, las raíces $\lambda_\pm=-\gamma\pm\iu\Omega$ recorren la circunferencia
$\abs{\lambda}=\omega$ (violeta) hasta chocar en $\lambda=-\omega$ para
$\nu=2$; después se separan sobre el eje real (naranja). Los puntos marcan
los tres casos de (b). (b) El régimen crítico,
$y=(1+x)\,\eul^{-x}$, es el regreso al equilibrio más rápido sin
oscilar.}
\label{fig:amortiguado}
\end{figure}

